% =========================================================================
% Cross-lingual Sense Transfer with Low-Resource Indic Languages:
% A Reproducible Kannada Word Sense Disambiguation Benchmark
% IEEE Access Extended Paper -- Target: 10-12 pages
% =========================================================================
\documentclass{ieeeaccess}
\usepackage{cite}
\usepackage{amsmath,amssymb,amsfonts}
\usepackage{algorithmic}
\usepackage{graphicx}
\usepackage{textcomp}
\usepackage{booktabs}
\usepackage{multirow}
\usepackage{array}
\usepackage{url}
\usepackage{longtable}

\def\BibTeX{{\rm B\kern-.05em{\sc i\kern-.025em b}\kern-.08em
    T\kern-.1667em\lower.7ex\hbox{E}\kern-.125emX}}

\begin{document}
\history{Date of publication June 2, 2026; date of current version June 2, 2026.}
\doi{10.1109/ACCESS.2026.0000000}

% =========================================================================
\title{Cross-lingual Sense Transfer with Low-Resource Indic Languages:
A Reproducible Kannada Word Sense Disambiguation Benchmark}

\author{%
\uppercase{Aditya Ranjan}\authorrefmark{1},
\uppercase{Arindam Gupta}\authorrefmark{1},
\uppercase{Gnanendra Naidu N}\authorrefmark{1},
\uppercase{Vijayalakshmi M.~N.}\authorrefmark{1}, and
\uppercase{S.~Anupama Kumar}\authorrefmark{1}}

\address[1]{Department of Artificial Intelligence and Machine Learning,
RV College of Engineering, Bengaluru 560059, India}

\tfootnote{Corresponding student authors: Aditya Ranjan
(adityaranjan.ai23@rvce.edu.in), Arindam Gupta
(arindamgupta.ai23@rvce.edu.in), and Gnanendra Naidu~N
(1gnanendranaidun.ai23@rvce.edu.in).
Faculty mentors: Dr.~Vijayalakshmi M.~N.\ and Dr.~S.~Anupama Kumar,
Associate Professors, Department of AIML.}

\markboth{Ranjan \headeretal: Cross-lingual Sense Transfer for Kannada WSD}
         {Ranjan \headeretal: Cross-lingual Sense Transfer for Kannada WSD}

\corresp{Corresponding author: Aditya Ranjan
(e-mail: adityaranjan.ai23@rvce.edu.in).}

% =========================================================================
\begin{abstract}
Word Sense Disambiguation (WSD) in low-resource languages remains an
unresolved bottleneck for high-quality natural language understanding.
Although multilingual transformer models have substantially improved
cross-lingual transfer, practical evidence for Kannada WSD under
reproducible, multi-task evaluation settings remains limited.
This paper presents a structured Kannada benchmark for cross-lingual sense
transfer and evaluates four multilingual encoders: XLM-R, mBERT, MuRIL,
and IndicBERT.
The resource package comprises a 45-word sense inventory with 150 sense
entries and 627 cleaned WiC-style sentence pairs over 44 polysemous target
words, covering nouns, verbs, and adjectives.
We evaluate three complementary formulations:
(i)~thresholded pair classification (WiC), using cosine similarity and
validation-tuned thresholds;
(ii)~unsupervised sense clustering scored with Adjusted Rand Index (ARI);
and
(iii)~gloss-based ranking scored with Mean Reciprocal Rank (MRR).

Results reveal clear task-dependent model behaviour: mBERT achieves the
strongest WiC classification (53.92\% test accuracy, 0.539 macro~F1);
MuRIL obtains the highest gloss-ranking performance (0.619 MRR); and
IndicBERT provides the best efficiency profile for constrained deployment
(33M parameters, fastest inference).
We provide qualitative t-SNE analyses for the highly polysemous target
\textit{kaala} across all four models, a structured error taxonomy, detailed
methodological rationale for target-token pooling over \texttt{[CLS]}
representations, and an energy-consumption comparison.
The findings demonstrate that robust Kannada WSD evaluation must jointly
measure pair discrimination, geometric separability, and lexical-gloss
alignment rather than relying on a single score.
All data, model outputs, figures, and result CSV files are released as
a reproducible artifact bundle.
\end{abstract}

\begin{keywords}
Cross-lingual transfer, Kannada NLP, low-resource language processing,
multilingual transformers, word sense disambiguation, WiC benchmark,
IndicBERT, MuRIL, XLM-R, mBERT.
\end{keywords}

\titlepgskip=-15pt
\maketitle

% =========================================================================
\section{Introduction}
\label{sec:intro}
% =========================================================================

Word Sense Disambiguation (WSD) determines which meaning of an ambiguous
lexical unit is active in a given context.
This decision influences downstream performance in machine translation,
semantic retrieval, dialogue systems, summarisation, and question answering.
English WSD has benefited from decades of lexical resource development
(WordNet, SemEval benchmarks) and large annotated corpora.
By contrast, Kannada---a Dravidian language with over 44 million speakers
and classical literary heritage---lacks sense-annotated resources and
standardised evaluation benchmarks.

Recent multilingual transformer models have created a practical pathway for
cross-lingual semantic transfer, but representation quality varies
substantially across tasks \cite{devlin2019,conneau2020}.
A model that is strong for pairwise semantic discrimination may not be
strongest for unsupervised cluster separation or gloss-alignment ranking.
Kannada WSD therefore requires evaluation protocols that capture
complementary semantic properties rather than reporting a single headline
metric.

The word \textit{kaala} illustrates the challenge: its five senses
(\textit{time/period}, \textit{leg}, \textit{death} (poetic),
\textit{black} (archaic), and \textit{season}) require both local
contextual cues and world-knowledge disambiguation.
Similarly, \textit{rasa} spans \textit{taste/flavour}, \textit{essence},
\textit{emotion}, \textit{liquid/nectar}, and \textit{artistic mood}---
a range that no surface-form heuristic can resolve.

\medskip\noindent This paper makes the following contributions:
\begin{itemize}
  \item A curated, low-resource Kannada WSD resource package with
        WiC-style pairs, sense inventory, example sentence inventory,
        and result traces.
  \item A cross-model comparison of XLM-R, mBERT, MuRIL, and IndicBERT
        under identical experimental constraints, covering three WSD
        formulations.
  \item Detailed multi-view evaluation using WiC classification, clustering
        ARI, and gloss MRR, with per-task analysis that demonstrates why
        single-metric reporting is misleading for low-resource WSD.
  \item Qualitative t-SNE geometry analysis for all four models on the
        polysemous target \textit{kaala}.
  \item An energy-consumption comparison motivating task-first model
        selection in resource-constrained Indic-language settings.
  \item A reproducible release structure containing data, outputs, and
        figures suitable for direct verification.
\end{itemize}

% =========================================================================
\section{Related Work}
\label{sec:related}
% =========================================================================

\subsection{Classical and Knowledge-Based WSD}

WSD research began with lexical-overlap heuristics, notably the Lesk
algorithm \cite{lesk1986}, which computed word overlap between a sentence
context and dictionary definitions.
The broad methodology landscape was formalised by Navigli \cite{navigli2009},
whose survey established that supervised approaches consistently outperform
knowledge-based ones given sufficient labelled data.
Standardised empirical comparison was later strengthened through unified
all-words evaluation protocols \cite{raganato2017}.

For Indic languages, the IndoWordNet project provided early sense inventories
for Hindi, Marathi, and Telugu, while Kannada coverage remained limited
\cite{kakwani2020}.
Kumar and Siddappa \cite{kumar2022} explored semantic-relation-based Kannada
WSD using hypernymy and hyponymy from IndoWordNet, but the approach relies on
lexical resources that are sparse for Kannada.

\subsection{Neural and Transformer-Based WSD}

Transformer pretraining \cite{devlin2019} transformed WSD by providing deep
contextual representations that naturally encode polysemy.
XLM-RoBERTa \cite{conneau2020} improved multilingual representation through
web-scale training, while MuRIL \cite{khanuja2021} specifically targeted
Indian languages by incorporating transliterated and romanised text.
IndicBERT \cite{kakwani2020} from AI4Bharat focused on a compact Indic-script
model.
Bevilacqua and Navigli \cite{bevilacqua2020} demonstrated that knowledge-graph
augmentation can push English WSD near 80\% F1---a ceiling this work uses as
contextual reference.
Singh et al.\ \cite{singh2021} evaluated multilingual transformers on Hindi
and Bengali WSD, finding model-task interaction effects similar to those we
report for Kannada.

\subsection{Low-Resource and Energy-Aware NLP}

Low-resource NLP has increasingly exploited cross-lingual transfer to avoid
annotation bottlenecks.
Energy-aware NLP \cite{strubell2019,schwartz2020} further highlights that
benchmark outcomes should be interpreted alongside efficiency and deployability
metrics---especially relevant for academic institutions in India without
large-scale GPU infrastructure.

% =========================================================================
\section{Dataset and Resource Construction}
\label{sec:dataset}
% =========================================================================

\subsection{Target Word Selection}

The target-word set was intentionally curated toward polysemous items with
three or more senses, covering 45 words: 18 nouns, 15 verbs, and 12 adjectives.
Words with only two senses were excluded to avoid inflated metrics under
threshold-based methods; higher-sense words better represent realistic
ambiguity.
Table~\ref{tab:wordstats} summarises the inventory by part of speech.

\begin{table}[t!]
\centering
\caption{Kannada Word Inventory by Part of Speech}
\label{tab:wordstats}
\begin{tabular}{l c l}
\toprule
\textbf{POS} & \textbf{Count} & \textbf{Example Words}\\
\midrule
Noun      & 18 & \textit{kaala} (5), \textit{rasa} (5), \textit{artha} (4)\\
Verb      & 15 & \textit{maadu} (5), \textit{tege} (4), \textit{kodu} (4)\\
Adjective & 12 & \textit{gaTTa} (3), \textit{mukta} (3), \textit{niila} (3)\\
\midrule
\textbf{Total} & \textbf{45} & Avg.\ 3.33 senses per word\\
\bottomrule
\end{tabular}
\end{table}

\subsection{Data Sources and Pair Construction}

The dataset pipeline was built from curated markdown and CSV resources derived
from IndoWordNet-oriented sense mappings and manually reviewed sentence
contexts.
Three primary resources were used:
(i)~a word inventory with sense definitions for all 45 words;
(ii)~an example sentence inventory with 5--10 annotated sentences per sense;
(iii)~a WiC-style pair file where pairs were constructed by combining same-sense
and different-sense sentence occurrences.

Pair generation followed controlled balancing: same-sense and different-sense
cases were both included to avoid skewed decision thresholds.
During cleaning, near-duplicate sentence forms were normalised to prevent
leakage through superficial text similarity.
This step is especially important in WiC-style evaluation where lexical overlap
can otherwise dominate contextual meaning.

\subsection{Pre-processing Pipeline}

The following pre-processing steps were applied:
\begin{enumerate}
  \item \textbf{Unicode normalisation:} NFC normalisation via Python's
        \texttt{unicodedata} module resolved encoding inconsistencies in
        Kannada script.
  \item \textbf{Whitespace normalisation:} Multiple spaces, tabs, and newlines
        were collapsed to single spaces.
  \item \textbf{Sentence key generation:} Trailing punctuation was stripped
        to enable duplicate detection across sources.
  \item \textbf{Gold sense assignment:} Sentences were matched against the
        inventory to propagate and verify gold-standard sense IDs.
  \item \textbf{Quality flagging:} Automated checks identified mixed-script
        sentences, short sentences ($<2$ tokens), and label conflicts.
\end{enumerate}

\subsection{Dataset Statistics and Quality Validation}

Table~\ref{tab:datasetstats} reports dataset-level statistics.

\begin{table}[t!]
\centering
\caption{Kannada WSD Benchmark Statistics}
\label{tab:datasetstats}
\begin{tabular}{lc}
\toprule
\textbf{Statistic} & \textbf{Value}\\
\midrule
Clean WiC pairs       & 627 \\
Raw WiC pairs         & 627 \\
Test split            & 497 \\
Validation split      & 126 \\
Unique target words   & 44 \\
Sense inventory words & 45 \\
Sense inventory rows  & 150 \\
Senses per word       & 3--5 (avg.\ 3.33)\\
\bottomrule
\end{tabular}
\end{table}

A systematic 20\% validation sample (every 5th pair) was extracted for
quality analysis.
Results showed 91\% correct labels, 9\% mislabelled pairs (mostly
adjective-sense blurring), and 5.6\% grammatically awkward sentences.
After correcting high-priority mislabels, projected quality improved to
90--92\%.
These figures indicate that dataset quality is the primary constraint on
absolute performance scores, consistent with the low-resource NLP literature.

% =========================================================================
\section{Methodology}
\label{sec:method}
% =========================================================================

\subsection{Model Suite}

We evaluate four multilingual encoders representing distinct pretraining
philosophies:
\begin{itemize}
  \item \textbf{mBERT} (\texttt{bert-base-multilingual-cased}, 110M params):
        broad multilingual baseline trained on 104 languages via masked
        language modelling on Wikipedia.
  \item \textbf{XLM-R} (\texttt{xlm-roberta-base}, 270M params):
        high-capacity web-scale multilingual model trained on 100 languages
        with 2.5~TB CommonCrawl data.
  \item \textbf{MuRIL} (\texttt{google/muril-base-cased}, 237M params):
        Indian-language-focused transformer incorporating 17 Indic languages,
        transliterated text, and romanised Indic text.
  \item \textbf{IndicBERT} (\texttt{ai4bharat/indic-bert}, 33M params):
        compact Indic-script model from AI4Bharat covering 12 major Indic
        languages.
\end{itemize}

\subsection{Target-Token Embedding Strategy}

Using target-token embeddings rather than global \texttt{[CLS]}-style vectors
is a deliberate methodological choice.
In many multilingual models, \texttt{[CLS]} emphasises sentence-level
semantics and may underweight local lexical ambiguity.
Target-token pooling better aligns the representation with WSD objectives,
particularly in morphologically rich Kannada text where the disambiguating
signal is concentrated near the ambiguous token.

For each sentence, the target polysemous word's character span is located
and all subword tokens overlapping that span are identified.
Their hidden-state vectors from the final transformer layer are averaged:
\[
\mathbf{x} = \frac{1}{|T_w|}\sum_{t \in T_w} \mathbf{h}_t
\]
where $T_w$ is the set of subword token indices overlapping the target word
span and $\mathbf{h}_t$ is the final-layer hidden state at position $t$.
The resulting vector is L2-normalised.

\subsection{Evaluation Protocol}

All models run over the same cleaned pair split with identical metric
definitions:
\begin{enumerate}
  \item \textbf{WiC classification:} For a pair $(x_i, x_j)$, compute
        \[
        s_{ij} = \frac{x_i^\top x_j}{\|x_i\|\|x_j\|}
        \]
        and predict $\hat{y}_{ij}=1$ if $s_{ij}\geq\tau$, else $0$.
        Threshold $\tau$ is tuned on the validation split to maximise macro~F1.
        Report validation accuracy, test accuracy, and test macro~F1.
  \item \textbf{Clustering:} Per-word K-means with $k$ equal to known sense
        cardinality.
        Cluster assignments are compared with gold sense labels using ARI.
        Report per-word ARI distribution.
  \item \textbf{Gloss ranking:} Each contextual embedding ranks all candidate
        gloss embeddings by cosine similarity.
        The reciprocal rank of the gold sense is accumulated; MRR is reported
        over all ranked instances.
\end{enumerate}

\subsection{Metric Complementarity}

Each metric captures a distinct failure mode:
\begin{itemize}
  \item \textbf{WiC accuracy/F1} is sensitive to decision threshold placement
        and class balance; it rewards clean bimodal similarity distributions.
  \item \textbf{ARI} is sensitive to geometric separability and cluster
        consistency across occurrences; it measures whether embedding space
        organises sense groups coherently.
  \item \textbf{MRR} is sensitive to ranking calibration and lexical alignment
        between contextual and gloss representations.
\end{itemize}
Reporting all three prevents overclaiming from one favorable metric and
supports transparent model selection for downstream use.

\subsection{Implementation Pipeline}

The pipeline is organised into five deterministic phases so that every
reported metric can be traced to intermediate artifacts:
\begin{description}
  \item[Phase 1:] Data parsing, Unicode normalisation, and pair-level record
    construction with stable identifiers.
  \item[Phase 2:] Batch embedding extraction using target-token pooling for
    all four models.
  \item[Phase 3:] WiC threshold selection on the validation split and
    test-set prediction export.
  \item[Phase 4:] Per-word K-means clustering and ARI computation.
  \item[Phase 5:] Gloss-candidate ranking, MRR aggregation, t-SNE
    visualisation, and report generation.
\end{description}

This staged design supports error localisation: data linkage errors appear
in Phase~1 outputs, representational drift in Phase~2 embeddings, and
metric-specific instability in later phases.
All intermediate CSV files and figures are retained in the \texttt{results/}
and \texttt{figures/} directories.

\subsection{Compute Environment}

All experiments were executed in CPU-compatible mode to ensure portability
in standard academic lab setups without GPU dependency.
This intentionally emphasises reproducibility over throughput and reflects
realistic constraints in low-resource institutional settings across India.

% =========================================================================
\section{Experimental Results}
\label{sec:results}
% =========================================================================

\subsection{WiC Classification Performance}

Table~\ref{tab:model_results_access} presents test-set WiC results for all
four models.

\begin{table}[t!]
\centering
\caption{WiC Classification Results (source: \texttt{results/model\_comparison.csv})}
\label{tab:model_results_access}
\begin{tabular}{lccccc}
\toprule
\textbf{Model} & \textbf{Thr.} & \textbf{Val Acc} & \textbf{Val F1}
  & \textbf{Test Acc} & \textbf{Test F1}\\
\midrule
mBERT      & 0.77 & 53.17\% & 0.532 & 53.92\% & 0.539\\
IndicBERT  & 0.90 & 51.59\% & 0.493 & 52.92\% & 0.489\\
MuRIL      & 0.10 & 46.83\% & 0.319 & 46.88\% & 0.319\\
XLM-R      & 0.85 & 46.83\% & 0.374 & 44.87\% & 0.338\\
\bottomrule
\end{tabular}
\end{table}

\begin{figure}[t!]
\centering
\includegraphics[width=0.84\textwidth]{figures/model_comparison.png}
\caption{Model-wise WiC performance on test accuracy and macro~F1. mBERT
achieves the strongest pairwise disambiguation behaviour in the current
benchmark; the gap to XLM-R and MuRIL is consistent across both metrics.}
\label{fig:model_access}
\end{figure}

mBERT achieves the highest combined test accuracy (53.92\%) and macro~F1
(0.539).
IndicBERT remains competitive despite its 33M-parameter compact footprint.
MuRIL's extremely low optimal threshold (0.10) indicates that its embedding
space produces uniformly compressed similarities for Kannada pairs---a
consequence of high-dimensional directional clustering rather than poor
representation quality.
XLM-R's lower WiC accuracy despite its large capacity (270M) suggests that
broad multilingual coverage does not directly translate to Kannada-specific
pairwise separability.

\subsection{Validation-to-Test Generalisation}

The validation-to-test performance gap is small for mBERT and IndicBERT
(less than 1.5 percentage points), indicating moderate generalisation under
current data scale.
XLM-R shows a slight drop from validation (46.83\%) to test (44.87\%),
suggesting that its threshold may be sensitive to small distributional
shifts.
These observations are consistent with cross-lingual embedding theory:
absolute similarity scales can differ across architectures, making a single
global threshold more suitable for some models than others.

\subsection{Clustering Quality (ARI)}

\begin{figure}[t!]
\centering
\includegraphics[width=0.84\textwidth]{figures/cluster_ari.png}
\caption{Per-word ARI distribution across all 44 target words for each model.
High variance reflects heterogeneous ambiguity difficulty; words with
semantically distant senses achieve higher ARI than semantically adjacent
senses.}
\label{fig:ari_access}
\end{figure}

Figure~\ref{fig:ari_access} shows that all models exhibit substantial
per-word ARI variance.
Some words (e.g., \textit{lobhi}, \textit{saadharana} for MuRIL and XLM-R)
achieve ARI~$=1.0$, indicating perfect cluster recovery.
Others (e.g., \textit{saadharana} for mBERT) achieve ARI~$=-0.5$, worse
than random assignment.
This heterogeneity reflects genuine variation in semantic-distance
structure across Kannada polysemous words rather than systematic model
failure, supporting the use of per-word diagnostics rather than aggregate
scores alone.

\subsection{Gloss Ranking Performance (MRR)}

\begin{table}[t!]
\centering
\caption{Gloss-Ranking MRR (source: \texttt{results/gloss\_metrics.csv})}
\label{tab:gloss_access}
\begin{tabular}{lcc}
\toprule
\textbf{Model} & \textbf{MRR} & \textbf{Ranked instances}\\
\midrule
MuRIL      & 0.619 & 294\\
mBERT      & 0.605 & 294\\
IndicBERT  & 0.592 & 294\\
XLM-R      & 0.570 & 294\\
\bottomrule
\end{tabular}
\end{table}

\begin{figure}[t!]
\centering
\includegraphics[width=0.84\textwidth]{figures/gloss_mrr.png}
\caption{Gloss-based MRR comparison across models. MuRIL leads, reversing
the WiC ranking and demonstrating that context-to-gloss alignment is a
distinct representational property from pairwise cosine separability.}
\label{fig:mrr_access}
\end{figure}

MuRIL's advantage on gloss alignment (0.619 MRR) reverses the WiC
classification ordering.
An MRR of 0.619 implies that the correct sense is on average ranked between
1st and 2nd position---a strong zero-shot result given that no sense-labelled
training data was used.
MuRIL's Indic-language pretraining and transliterated-text augmentation
appear to produce contextual vectors that align more directly with
dictionary-style gloss language than the representations from general
multilingual encoders.

\subsection{Cross-Metric Model Ranking Summary}

Table~\ref{tab:rankings} summarises model rankings across all three tasks.

\begin{table}[t!]
\centering
\caption{Task-Dependent Model Rankings}
\label{tab:rankings}
\begin{tabular}{lccc}
\toprule
\textbf{Rank} & \textbf{WiC (Acc/F1)} & \textbf{Clustering (ARI)} & \textbf{Gloss (MRR)}\\
\midrule
1st & mBERT     & XLM-R / MuRIL & MuRIL\\
2nd & IndicBERT & MuRIL / XLM-R & mBERT\\
3rd & MuRIL     & mBERT         & IndicBERT\\
4th & XLM-R     & IndicBERT     & XLM-R\\
\bottomrule
\end{tabular}
\end{table}

The cross-task rank shifts confirm that no single encoder dominates all
semantic evaluation views.
This finding directly challenges the common practice of reporting a single
WiC accuracy as the definitive Kannada WSD performance score.

\subsection{t-SNE Qualitative Geometry}

\begin{figure*}[t!]
\centering
\includegraphics[width=0.49\textwidth]{figures/tsne_indicbert_kala.png}
\includegraphics[width=0.49\textwidth]{figures/tsne_mbert_kala.png}
\includegraphics[width=0.49\textwidth]{figures/tsne_muril_kala.png}
\includegraphics[width=0.49\textwidth]{figures/tsne_xlmr_kala.png}
\caption{t-SNE projections for the highly polysemous target \textit{kaala}
(5 senses: time, leg, death, black, season) across all four models.
Differences in cluster compactness, overlap, and boundary sharpness provide
qualitative explanations for the ARI and MRR differences reported in
Tables~\ref{tab:gloss_access} and~\ref{tab:rankings}. Top-left: IndicBERT;
top-right: mBERT; bottom-left: MuRIL; bottom-right: XLM-R.}
\label{fig:tsne_access}
\end{figure*}

The t-SNE plots (Figure~\ref{fig:tsne_access}) reveal that cluster compactness
and separability vary substantially by architecture.
mBERT produces relatively sharp boundaries for semantically distant senses
(time vs.\ leg), explaining its WiC advantage.
MuRIL shows tighter local groupings near sense centroids, consistent with its
stronger gloss-alignment behaviour.
IndicBERT displays the most compact overall layout, reflecting its compact
embedding space; XLM-R shows broader scatter, potentially due to over-generalised
multilingual representations.
Across all models, the time/season sense boundary is the least separable,
confirming that semantically adjacent senses represent the primary difficulty
for Kannada WSD.

% =========================================================================
\section{Extended Analysis}
\label{sec:analysis}
% =========================================================================

\subsection{Why mBERT Leads WiC Classification}

mBERT's WiC advantage stems from a cleaner bimodal distribution of cosine
similarity scores for same-sense vs.\ different-sense pairs at the tuned
threshold of~0.77.
mBERT's training on 104 languages via masked language modelling on Wikipedia
produces a moderately diverse but structured similarity space where Kannada
contextual tokens occupy distinct regions for different senses.
XLM-R, despite its larger capacity, produces a more compressed similarity
distribution where a single global threshold is less effective.

\subsection{Why MuRIL Leads Gloss Ranking}

MuRIL's stronger MRR reflects better alignment between contextual sentence
vectors and gloss definition vectors for Indic text.
The model was trained with explicit augmentation for transliterated and
romanised Indic text, which increases overlap between dictionary-style gloss
language (often written in a formal register) and naturally occurring sentence
contexts.
Furthermore, MuRIL's Kannada-specific tokenisation---with fewer subword
splits than mBERT for morphologically rich Kannada forms---produces
target-token embeddings that better reflect the word's semantic role.

\subsection{ARI Distribution Interpretation}

ARI variance across target words indicates heterogeneous ambiguity difficulty.
Words with semantically distant senses (e.g., \textit{chakra}: wheel vs.\
administrative district; ARI up to 0.60 for XLM-R) form cleaner clusters.
Words with closely related senses (e.g., \textit{maadu}: to do vs.\ to build;
ARI near 0 across all models) collapse into overlapping manifolds.
This behavior is expected from the cognitive science of polysemy: senses that
developed from common etymological roots through metonymy or metaphor are
harder to separate than senses that are fully unrelated.

\subsection{Energy and Efficiency Analysis}

Table~\ref{tab:energy} summarises the computational profile of each model.

\begin{table}[t!]
\centering
\caption{Computational Profile During Embedding Extraction (CPU)}
\label{tab:energy}
\begin{tabular}{lccl}
\toprule
\textbf{Model} & \textbf{Params} & \textbf{Rel.\ Speed} & \textbf{Est.\ Energy}\\
\midrule
IndicBERT & 33M  & 1.0$\times$ (fastest) & $\approx$0.8 kWh / 10k sent.\\
mBERT     & 110M & 1.8$\times$           & $\approx$1.5 kWh / 10k sent.\\
MuRIL     & 237M & 2.5$\times$           & $\approx$2.1 kWh / 10k sent.\\
XLM-R     & 270M & 3.0$\times$ (slowest) & $\approx$3.4 kWh / 10k sent.\\
\bottomrule
\end{tabular}
\end{table}

IndicBERT is the most energy-efficient model (33M parameters, fastest CPU
inference).
mBERT offers the best accuracy-to-energy ratio: it achieves the highest WiC
scores at only 1.8$\times$ the IndicBERT cost.
MuRIL's Indian-language specialisation justifies its additional compute for
gloss-ranking tasks.
XLM-R is the most expensive model and---counter-intuitively---delivers the
weakest Kannada WSD results, suggesting that broad multilingual capacity does
not substitute for language-specific pretraining in low-resource settings
\cite{strubell2019,schwartz2020}.

% =========================================================================
\section{Error Analysis}
\label{sec:error}
% =========================================================================

\subsection{Error Taxonomy}

Systematic inspection of mismatch cases revealed three primary failure
categories:

\subsubsection{Sense Proximity Errors}
Semantically adjacent senses (e.g., \textit{time} vs.\ \textit{season} for
\textit{kaala}, or \textit{doing} vs.\ \textit{making} for \textit{maadu})
produce overlapping embedding regions even for sentences that human annotators
rated as unambiguous.
These cases account for the majority of WiC false negatives across all models.

\subsubsection{Tokeniser Fragmentation Errors}
Kannada's agglutinative morphology produces long subword chains in all four
tokenisers.
Suffixes encoding tense, case, and aspect may receive separate subword tokens,
reducing the fidelity of target-token pooling when the disambiguating suffix
is morphologically distant from the base form.
Models with fewer average Kannada subword tokens per word (MuRIL, IndicBERT)
tend to exhibit lower fragmentation error rates.

\subsubsection{Context Sparsity Errors}
Short or syntactically minimal sentences (fewer than 6 tokens) provide
insufficient disambiguating signal.
The surrounding lexical context is often neutral (function words, common verbs)
and does not strongly activate sense-specific features.
These cases show near-random cosine similarity scores across all models.

\subsection{POS-Level Patterns}

Verbs generally exhibit lower WiC accuracy than nouns, consistent with
their stronger dependence on argument structure and discourse context.
Adjectives show moderate performance, with errors concentrated in
gradable senses (e.g., \textit{gaTTa}: hard vs.\ rigid vs.\ strong).
This POS-level pattern suggests that future Kannada WSD systems should
incorporate part-of-speech-aware threshold calibration.

% =========================================================================
\section{Limitations and Future Directions}
\label{sec:limits}
% =========================================================================

\subsection{Current Limitations}

Three limitations constrain the current benchmark:
\begin{enumerate}
  \item \textbf{Moderate corpus scale:} 627 pairs and 44 target words provide
        a reasonable starting point but insufficient statistical power for
        confident cross-model significance testing.
  \item \textbf{Gloss quality variability:} Gloss richness differs across
        senses, biasing MRR toward frequent senses with more descriptive
        glosses.
  \item \textbf{Context locality:} The pipeline uses sentence-level contexts
        without discourse-level cues, which limits disambiguation for words
        whose senses are determined by broader pragmatic context.
\end{enumerate}

\subsection{Future Directions}

Immediate extensions include:
\begin{itemize}
  \item \textbf{Supervised fine-tuning:} Adapting MuRIL or mBERT on the
        curated Kannada WiC training set, following the few-shot WSD paradigm.
  \item \textbf{Lexical inventory expansion:} Broadening target coverage to
        domain-specialised vocabulary (legal, medical, agricultural).
  \item \textbf{Cross-Indic transfer:} Exploiting high-resource Telugu or
        Tamil WSD data to improve rare-sense handling for Kannada through
        related-language transfer.
  \item \textbf{Calibrated confidence estimation:} Implementing temperature
        scaling or Platt calibration to produce reliable uncertainty estimates
        for deployment.
  \item \textbf{Discourse context integration:} Extending the pipeline to
        multi-sentence contexts for disambiguation of pragmatically sensitive
        senses.
\end{itemize}

% =========================================================================
\section{Practical Deployment Guidelines}
\label{sec:deploy}
% =========================================================================

For production and near-production settings, model selection should follow
task-first constraints:
\begin{itemize}
  \item \textbf{Binary sense decision:} Use mBERT with a validation-tuned
        threshold.
        It delivers the strongest WiC accuracy and the best
        accuracy-to-compute ratio among the four models.
  \item \textbf{Glossary or dictionary integration:} Use MuRIL for gloss-based
        ranking or dictionary-assisted interfaces.
        Its Indic-language alignment provides the strongest context-to-gloss
        coupling.
  \item \textbf{Compute-constrained deployment:} Use IndicBERT with conservative
        confidence thresholds.
        Its 33M-parameter footprint and 3$\times$ speed advantage make it
        practical for mobile or edge inference.
\end{itemize}

In all cases, periodic error audits over newly collected domain sentences are
recommended to detect silent semantic drift as language usage evolves.

% =========================================================================
\section{Reproducibility Statement}
\label{sec:repro}
% =========================================================================

This submission package includes all required source artifacts to reproduce
every table and figure:
\begin{itemize}
  \item \texttt{data/}: cleaned and raw pair datasets, sense inventory, example
        sentence inventory, and validation sample.
  \item \texttt{results/}: model-level summaries (\texttt{model\_comparison.csv},
        \texttt{gloss\_metrics.csv}) and pair-level traces
        (\texttt{pair\_level\_predictions.csv},
        \texttt{clustering\_summary.csv}).
  \item \texttt{figures/}: chart outputs (model comparison, cluster ARI, gloss
        MRR) and t-SNE visualisations for all four models.
\end{itemize}

All manuscript figure paths are relative (e.g., \texttt{figures/...}) so that
Overleaf and local builds remain consistent without machine-specific
absolute paths.
The pipeline is implemented in a single reproducible script
(\texttt{run\_pipeline.py}) that can be executed with:
\begin{verbatim}
  python run_pipeline.py --phase all
\end{verbatim}

% =========================================================================
\section{Conclusion}
\label{sec:conc}
% =========================================================================

This study establishes a reproducible Kannada cross-lingual WSD benchmark
and demonstrates that model quality is evaluation-task dependent.
mBERT leads WiC pair discrimination; MuRIL leads gloss-based ranking;
IndicBERT offers the strongest deployment efficiency.
No single model dominates all three evaluation views.

The broader implication is methodological: low-resource semantic evaluation
benefits from multi-perspective measurement.
Pair classification rewards threshold separability; clustering rewards
geometric coherence; gloss ranking rewards lexical-semantic alignment.
A robust Kannada WSD leaderboard should jointly report all three.

The released artifact bundle is intended to support follow-up research on
Kannada and to serve as a practical template for other Indic languages where
annotation budgets are limited.
Future work will prioritise supervised fine-tuning, richer sense examples,
and cross-Indic few-shot transfer to close the gap between zero-shot
performance and practical deployment requirements.

% =========================================================================
\begin{thebibliography}{99}

\bibitem{lesk1986}
M. Lesk, ``Automatic sense disambiguation using machine readable dictionaries:
How to tell a pine cone from an ice cream cone,'' in
\emph{Proc. 5th Annu. Int. Conf. Syst. Documentation}, 1986,
pp.~24--26, doi:\,10.1145/318723.318728.

\bibitem{navigli2009}
R. Navigli, ``Word sense disambiguation: A survey,''
\emph{ACM Comput. Surv.}, vol.~41, no.~2, pp.~1--69, 2009,
doi:\,10.1145/1459352.1459355.

\bibitem{raganato2017}
A. Raganato, J. Camacho-Collados, and R. Navigli, ``Word sense disambiguation:
A unified evaluation framework and empirical comparison,'' in
\emph{Proc. EACL}, 2017, pp.~99--110, doi:\,10.18653/v1/E17-1010.

\bibitem{devlin2019}
J. Devlin, M.-W. Chang, K. Lee, and K. Toutanova, ``BERT: Pre-training of
deep bidirectional transformers for language understanding,'' in
\emph{Proc. NAACL-HLT}, 2019, pp.~4171--4186, doi:\,10.18653/v1/N19-1423.

\bibitem{conneau2020}
A. Conneau, K. Khandelwal, N. Goyal, V. Chaudhary, G. Wenzek,
F. Guzm\'{a}n, E. Grave, M. Ott, L. Zettlemoyer, and V. Stoyanov,
``Unsupervised cross-lingual representation learning at scale,'' in
\emph{Proc. ACL}, 2020, pp.~8440--8451, doi:\,10.18653/v1/2020.acl-main.747.

\bibitem{kakwani2020}
D. Kakwani, A. Kunchukuttan, S. Golla, G. N. C., A. Bhattacharyya,
M.~M. Khapra, and P. Kumar, ``IndicNLPSuite: Monolingual corpora, evaluation
benchmarks and pre-trained multilingual language models for Indian languages,''
in \emph{Findings of EMNLP}, 2020, pp.~4948--4961,
doi:\,10.18653/v1/2020.findings-emnlp.445.

\bibitem{khanuja2021}
S. Khanuja, D. Bansal, S. Mehtani, S. Khosla, A. Dey, B. Gopalan,
D.~K. Margam, P. Aggarwal, R.~T. Nagipogu, S. Dave, S. Gupta,
S.~C.~B. Gali, V. Subramanian, and P. Talukdar, ``MuRIL: Multilingual
representations for Indian languages,'' \emph{arXiv preprint
arXiv:2103.10730}, 2021.

\bibitem{bevilacqua2020}
M. Bevilacqua and R. Navigli, ``Breaking through the 80\% glass ceiling:
Raising the state of the art in word sense disambiguation by incorporating
knowledge graph information,'' in \emph{Proc. ACL}, 2020, pp.~2854--2864,
doi:\,10.18653/v1/2020.acl-main.255.

\bibitem{kumar2022}
B.~H.~M. Kumar and M. Siddappa, ``Kannada word sense disambiguation using
semantic relations,'' \emph{Int. J. Comput. Appl.}, vol.~183, no.~48,
pp.~12--19, 2022, doi:\,10.5120/ijca2022921887.

\bibitem{strubell2019}
E. Strubell, A. Ganesh, and A. McCallum, ``Energy and policy considerations
for deep learning in NLP,'' in \emph{Proc. ACL}, 2019, pp.~3645--3650,
doi:\,10.18653/v1/P19-1355.

\bibitem{schwartz2020}
R. Schwartz, J. Dodge, N.~A. Smith, and O. Etzioni, ``Green AI,''
\emph{Commun. ACM}, vol.~63, no.~12, pp.~54--63, 2020,
doi:\,10.1145/3381831.

\bibitem{singh2021}
A. Singh, S. Husain, and M. Shrivastava, ``Multilingual transformers for
Indian language word sense disambiguation,'' in
\emph{Proc. ACL Student Research Workshop}, 2021, pp.~74--83,
doi:\,10.18653/v1/2021.acl-srw.10.

\end{thebibliography}

\EOD
\end{document}
